\section{Transformações Lineares}

O mais simples tipo de transformação que pode ser aplicado a um atributo $X$ quantitativo
é uma transformação linear na forma $aX+b$, com $a, b \in \real$.
Ela pode ser interpretada como um escalonamento seguido de um deslocamento.
Os valores dos parâmetros $a$ e $b$ dependem da finalidade da exploração.

Analisemos o impacto de uma transformação linear em algumas medidas de resumo.
No Teorema \ref{thm:propriedades esperanca} vimos que a esperança é um operador linear.
Como a média é a esperança de $X$ sob a probabilidade empírica, isso implica que ela também
o seja; ainda assim, a verificação direta é imediata.

\begin{prop:media linear}
\label{prop:media linear}
Se $Y=aX+b$ é uma transformação aplicada a um atributo $X$, então a média de $Y$ é
$a \bar X + b$.
\end{prop:media linear}
\begin{proof}
Seja $n$ o número de linhas do quadro de $X$.
\[
    \frac 1 n \sum_{i=1}^n (ax_i + b) = \frac a n \sum_{i=1}^n x_i + b = a \bar X + b
\]
\end{proof}

Algo análogo à transformação linear da média ocorre com a mediana.

\begin{prop:mediana linear}
\label{prop:mediana linear}
Se $Y=aX+b$ é uma transformação aplicada a um atributo $X$, então a mediana de $Y$ é
$a \cdot \med X + b$.
\end{prop:mediana linear}
\begin{proof}
Se $a \ge 0$, então a transformação $x \mapsto ax + b$ é não decrescente, e a ordenação
das estatísticas de ordem não é alterada, isto é, $\eo{y}{k} = a \eo{x}{k} + b$ para todo
$k \in [n]$.
Sendo $n$ o número de linhas do quadro de $X$, se $n$ é ímpar, então a mediana é
\[
    a x_{(\frac{n+1}{2})} + b.
\]
Caso seja par,
\[
   \frac{(a \eo{x}{\frac n 2} + b) + (a \eo{x}{\frac{n}{2}+1} + b)}{2} =
    a\frac{(\eo{x}{\frac n 2} + \eo{x}{\frac{n}{2}+1})}{2} + b =
    a \cdot \med{X} + b.
\]
Quando $a < 0$, ocorre uma inversão na ordem das estatísticas de ordem, de modo que
$\eo{y}{k} = a \eo{x}{n+1-k} + b$ para todo $k \in [n]$.
Todavia, a troca de índices $k \mapsto n+1-k$ leva $\frac{n+1}{2}$ nele mesmo e permuta
$\frac n 2$ com $\frac n 2 + 1$.
Assim, se $n$ é ímpar, $\eo{y}{\frac{n+1}{2}} = a \eo{x}{\frac{n+1}{2}} + b$, e, se $n$
é par, as estatísticas usadas para computar a mediana de $Y$ são as mesmas do caso anterior
em ordem trocada, o que não altera a sua média aritmética.
Em ambos os casos, a mediana de $Y$ é $a \cdot \med X + b$.
\end{proof}

A moda acompanha a transformação de forma ainda mais direta, pois depende apenas das
frequências dos valores, e não de sua ordenação.

\begin{prop:moda linear}
\label{prop:moda linear}
Se $Y=aX+b$ é uma transformação aplicada a um atributo $X$ e $m$ é uma moda amostral de
$X$, então $am + b$ é uma moda amostral de $Y$.
\end{prop:moda linear}
\begin{proof}
Se $a = 0$, então $Y$ é constante e igual a $b = am + b$, que é, pois, sua única moda.
Se $a \neq 0$, então a função $x \mapsto ax + b$ é injetora, de modo que $Y$ assume o
valor $ax + b$ exatamente nas linhas em que $X$ assume o valor $x$, isto é,
$N_Y(\{ax+b\}) = N_X(\{x\})$ para todo $x \in X(D)$.
Como $m$ maximiza $N_X(\{x\})$, segue que $am + b$ maximiza $N_Y(\{y\})$.
\end{proof}
Quando $a \neq 0$, a correspondência é uma bijeção entre as modas de $X$ e as de $Y$.

A mediana é um caso particular dos quantis empíricos, e o comportamento visto acima se
estende a eles, com a ressalva de que, para $a < 0$, a inversão da ordem troca o quantil
de probabilidade $p$ pelo de probabilidade $1 - p$.
Denotaremos por $q_p(X)$ o quantil empírico de probabilidade $p$ de $X$ obtido pelo
algoritmo \ref{alg:quantil}.

\begin{prop:quantil linear}
\label{prop:quantil linear}
Seja $Y=aX+b$ uma transformação aplicada a um atributo $X$ de um quadro com $n \ge 2$
linhas, e seja $p \in [0, 1]$.
Se $a \ge 0$, então $q_p(Y) = a\, q_p(X) + b$; se $a < 0$, então
$q_p(Y) = a\, q_{1-p}(X) + b$.
\end{prop:quantil linear}
\begin{proof}
Seja $\ell_X : [1, n] \to \real$ a função linear por partes que liga os pontos
$(k, \eo{x}{k})$, para $k \in [n]$, isto é,
\[
    \ell_X(t) = (1 - (t - k))\, \eo{x}{k} + (t - k)\, \eo{x}{k+1},
    \quad t \in [k, k+1].
\]
Pelo algoritmo \ref{alg:quantil}, com $h = (n-1)p + 1$, temos $q_p(X) = \ell_X(h)$,
inclusive no caso $p = 1$, em que $q_1(X) = \eo{x}{n} = \ell_X(n)$.
Defina $\ell_Y$ de forma análoga para $Y$.

Se $a \ge 0$, então $\eo{y}{k} = a\eo{x}{k} + b$, como visto na prova da proposição
\ref{prop:mediana linear}.
As funções $\ell_Y$ e $a\ell_X + b$ são lineares em cada intervalo $[k, k+1]$ e coincidem
nos inteiros, logo são iguais, e $q_p(Y) = \ell_Y(h) = a\, q_p(X) + b$.

Se $a < 0$, então $\eo{y}{k} = a\eo{x}{n+1-k} + b$.
A função $t \mapsto a\,\ell_X(n+1-t) + b$ também é linear em cada intervalo $[k, k+1]$,
pois $n + 1 - t$ percorre o intervalo $[n-k, n+1-k]$, e vale $a\eo{x}{n+1-k} + b = \eo{y}{k}$
em cada inteiro $k$, de modo que $\ell_Y(t) = a\,\ell_X(n+1-t) + b$.
Como $n + 1 - h = (n-1)(1-p) + 1$, segue que
\[
    q_p(Y) = a\,\ell_X((n-1)(1-p) + 1) + b = a\, q_{1-p}(X) + b.
\]
\end{proof}
Para $p = 0.5$, temos $1 - p = p$, e a proposição \ref{prop:quantil linear} recupera a
proposição \ref{prop:mediana linear}.

A transformação linear da variância amostral é análoga à mostrada para aquela do
Teorema \ref{prop:propriedades variancia}.

\begin{prop:variancia linear}
\label{prop:variancia linear}
Se $Y=aX+b$ é uma transformação aplicada a um atributo $X$, então $s_Y^2=a^2 s_X^2$.
\end{prop:variancia linear}
\begin{proof}
Se $n$ é o número de linhas do quadro de $X$, então
\[
   \frac{1}{n-1} \sum_{i=1}^n (ax_i+b - (a\bar X+b))^2 =
   \frac{a^2}{n-1} \sum_{i=1}^n (x_i - \bar X)^2 =
   a^2 s_X^2
\]
\end{proof}

O desvio absoluto médio, o intervalo interquartílico e o desvio absoluto mediano
comportam-se da mesma maneira que o desvio-padrão:
são multiplicadas por $|a|$ e não dependem do deslocamento $b$.
Isso é esperado, pois um deslocamento move todos os dados juntos, sem alterar as
distâncias entre eles, ao passo que o escalonamento multiplica todas essas distâncias
por $|a|$.

\begin{prop:dispersao linear}
\label{prop:dispersao linear}
Seja $Y=aX+b$ uma transformação aplicada a um atributo $X$ de um quadro com $n \ge 2$
linhas.
Então
\begin{gather*}
    \dm{Y} = |a| \dm{X}, \quad s_Y = |a|\, s_X, \\
    d_q(Y) = |a|\, d_q(X) \quad \text{e} \quad
    \mathrm{MAD}(Y) = |a|\, \mathrm{MAD}(X),
\end{gather*}
onde $d_q(X) = q_{0.75}(X) - q_{0.25}(X)$ é o intervalo interquartílico empírico.
\end{prop:dispersao linear}
\begin{proof}
Pela proposição \ref{prop:media linear}, $y_i - \bar Y = a(x_i - \bar X)$, logo
\[
    \dm{Y} = \frac 1 n \sum_{i=1}^n |a|\,|x_i - \bar X| = |a| \dm{X}.
\]
A igualdade $s_Y = |a|\, s_X$ segue da proposição \ref{prop:variancia linear} tomando a
raiz quadrada.
Para o intervalo interquartílico, se $a \ge 0$, a proposição \ref{prop:quantil linear} dá
\[
    d_q(Y) = (a\, q_{0.75}(X) + b) - (a\, q_{0.25}(X) + b) = a\, d_q(X),
\]
e, se $a < 0$, ela dá
\[
    d_q(Y) = (a\, q_{0.25}(X) + b) - (a\, q_{0.75}(X) + b) = -a\, d_q(X) = |a|\, d_q(X).
\]
Por fim, pela proposição \ref{prop:mediana linear}, $|y_i - \med Y| = |a|\,|x_i - \med X|$.
Assim, os desvios absolutos de $Y$ em relação à sua mediana são os de $X$ multiplicados
por $|a| \ge 0$, e, aplicando novamente a proposição \ref{prop:mediana linear}, com
coeficiente $|a|$ e deslocamento $0$, obtemos $\mathrm{MAD}(Y) = |a|\, \mathrm{MAD}(X)$.
\end{proof}

Como consequência, o coeficiente de variação satisfaz $\mathrm{CV}(aX) = \mathrm{CV}(X)$
para $a > 0$, mas não é preservado por deslocamentos, já que $s_X$ não depende de $b$,
enquanto a média depende, como já observado na seção de medidas de dispersão.

Por fim, os coeficientes de assimetria não dependem nem da escala, nem da posição dos
dados, mas apenas do sentido da transformação.

\begin{prop:assimetria linear}
\label{prop:assimetria linear}
Seja $Y=aX+b$, com $a \neq 0$, uma transformação aplicada a um atributo $X$, e seja $c$ um
dos coeficientes de assimetria de Pearson I, Pearson II, quartílico ou de Fisher, definido
para $X$.
Então $c$ está definido para $Y$, e
\[
    c(Y) = \frac{a}{|a|}\, c(X).
\]
\end{prop:assimetria linear}
\begin{proof}
Pelas proposições \ref{prop:media linear}, \ref{prop:mediana linear},
\ref{prop:moda linear} e \ref{prop:dispersao linear}, as diferenças $\bar Y - \mo Y$ e
$\bar Y - \med Y$ são iguais às respectivas diferenças de $X$ multiplicadas por $a$, ao
passo que $s_Y = |a|\, s_X$, o que prova o resultado para os coeficientes de Pearson,
com $s_Y > 0$ sempre que $s_X > 0$.

Para o coeficiente quartílico, se $a > 0$, a proposição \ref{prop:quantil linear} dá
\[
    q_{0.25}(Y) + q_{0.75}(Y) - 2\med Y
        = a\left(q_{0.25}(X) + q_{0.75}(X) - 2\med X\right);
\]
se $a < 0$, os quartis trocam de papel, mas a soma $q_{0.25} + q_{0.75}$ é a mesma, e a
igualdade acima continua válida.
Como $d_q(Y) = |a|\, d_q(X)$, o resultado segue.

Para o coeficiente de Fisher, temos $(y_i - \bar Y)^3 = a^3 (x_i - \bar X)^3$, e, sendo
$\sigma_X$ e $\sigma_Y$ os desvios-padrão com divisor $n$ de $X$ e $Y$, pelo mesmo
argumento da proposição \ref{prop:variancia linear}, $\sigma_Y = |a|\,\sigma_X > 0$.
O fator de correção depende apenas de $n$, logo o coeficiente de $Y$ é o de $X$
multiplicado por $\frac{a^3}{|a|^3} = \frac{a}{|a|}$.
\end{proof}

Em particular, para $a > 0$, nenhum desses coeficientes se altera, o que formaliza a
afirmação de que transformações lineares não mudam a forma da distribuição.
Para $a < 0$, a distribuição é refletida, e a cauda mais longa muda de lado.

Dentre as transformações lineares mais usadas em análise exploratória está a padronização,
por vezes também chamada de normalização.
Ela é usada para ajustar os valores de uma sequência de atributos a uma escala comum
para facilitar a comparação entre eles.

Um exemplo de padronização é a transformação $(X - \bar X)/s_X$, definida quando
$s_X > 0$, isto é, quando $X$ não é constante.
Ela produz um novo atributo com média $0$ e desvio-padrão amostral igual a $1$, o que
segue das Proposições \ref{prop:media linear} e \ref{prop:variancia linear} com
$a = 1/s_X$ e $b = -\bar X/s_X$, de forma análoga ao exposto na Proposição
\ref{prop:padronizacao}.
Essa transformação é adequada quando o atributo se aproxima de uma distribuição normal,
de modo que o atributo resultante aproximar-se-á de uma distribuição normal padrão.
Também é útil quando aplicamos modelos sensíveis à escala dos atributos, como aqueles
baseados em distâncias.
Note, porém, que, por ser linear com $a = 1/s_X > 0$, a padronização não altera a forma da
distribuição, conforme a proposição \ref{prop:assimetria linear}, de modo que não torna
normal um atributo que não o seja.
Além disso, como $\bar X$ e $s_X$ não são robustos, a presença de \textit{outliers} infla
$s_X$ e comprime os valores típicos num intervalo estreito, prejudicando a comparação entre
atributos.

Uma variação robusta da padronização utiliza a mediana e o IQR no lugar da média e do
desvio-padrão, respectivamente, dada por $(X - \med X)/d_q$, definida quando $d_q > 0$.
Essa condição pode falhar mesmo para atributos não constantes: para $X = (0, 0, 0, 0, 1)$,
o método de interpolação dá $q_{0.25}(X) = q_{0.75}(X) = 0$.
Tomando $a = 1/d_q$ e $b = -\med X / d_q$, o resultado é um atributo com mediana $0$,
pela proposição \ref{prop:mediana linear}, e IQR igual a $1$, pela proposição
\ref{prop:dispersao linear}.
Sua aplicação é recomendada quando os dados possuem \textit{outliers} que não podem ser
descartados, possuem distribuição assimétrica ou cauda pesada.

Outro exemplo é a normalização min-max, que escala os valores para o intervalo $[0,1]$,
definida quando $\max X > \min X$ e dada por
\[
    \frac{X - \min X}{\max X - \min X}.
\]
Sua vantagem sobre a padronização tradicional é que os dados são ajustados para um
intervalo fixo e não negativo, porém também sofre na presença de \textit{outliers}.
